% This is the document class for AEBRs. Note that this format is interchangable with
% report. Any AEBR will compile as a dragonfly-report and vice versa.
% The draft option will also work. 
\documentclass{aebr}

% The AEBR requires two titles. The first is the short title which goes on the 
% page footers. The second is the long title which goes on the cover page. 
% You are only allowed to use the the LaTeX macros \\ and \emph inside of titles
% safely. Others may result in unpredictable behaviour. 
\title{Example}{An Example\\ Report about \emph{sheep} but that also has such
a long title that it actually has to wrap!}

% Each AEBR has a subtitle
\subtitle{A report for testing purposes}

%The date must be split into two parts for compatability with reports. 
% The year is specified with \date, and the month is specified with \reportmonth
\date{2014}
\reportmonth{May}

% Authors must be listed with their full names separated by \and.
\author{John Smith \and Jane E. Smith}

% The subcaption package is used to create subfigures in this example.
\usepackage{subcaption}

% The lipsum package is used to generate filler text.
\usepackage{lipsum}

% Each report has an ISSN, ISBN and report number. If any of these are not 
% included a place holder value will be used.
\issn{1234-5678}
\isbn{12-34567-89}
\reportno{79}

% This specified the location of the bibliography file
\addbibresource{test.bib}

% The document begins here
\begin{document}

% Generate the title page
\maketitle

% Generate the table of contents. This should be done for AEBRs.
\tableofcontents

% An AEBR must contain an executive summary. Which \summary will generate a 
% title and contents entry for. 
\summary

% The summary must start with a citation for the current document which can be 
% automatically generated using \citeself.
\citeself

This has some boring text in it



\section{The First}

%filler text
\lipsum[1]


\begin{figure}[h]
  \includegraphics[width=40mm,height=40mm]{pattern}
  \caption{This is a caption that is longer than one line. It is improtant to test 
  this behavoir as this a common feature of the dragon fly templates.}
\end{figure}

\subsection{Details}

\lipsum[2]

% An example of how to do subfigures using subcaption.
\begin{figure}[h]

  \begin{subfigure}{0.45\textwidth}
  \includegraphics[width=40mm,height=40mm]{pattern}
  \caption{This is a caption that is longer than one line. It is improtant to test 
  this behavoir as this a common feature of the dragon fly templates.}
\end{subfigure}\qquad
  \begin{subfigure}{0.45\textwidth}
  \includegraphics[width=40mm,height=40mm]{pattern}
  \caption{This is a caption that is longer than one line. It is improtant to test 
  this behavoir as this a common feature of the dragon fly templates.}
\end{subfigure}
\end{figure}

\subsubsection{What's This?}
\lipsum[3] 
\paragraph{A minor heading}
\lipsum[5] 
\subparagraph{A really minor heading}
\lipsum[4]

Apa comes with a number of ways of citing things.
The interesting ones are:

% Examples of how to do citations. Note the cite, citet and citep work as standard. 
\begin{itemize}
  \item \cite{abraham_summary_98-09}
  \item \citet{baker_nzclassification_2010}
  \item \citep{doc_sealion_2009}
  \item \parencite{gales_phocarctos_2008}
  \item \nptextcite{ mpi_review_2012}
  \item \fullcite{mpi_review_2012}
  \item \fullcitebib{mpi_review_2012}
  \item \citeyear{robertson_population_2011}
  %\item \footcite{mpi_review_2012}
  \item \textcite{roe_necropsy_2007}
  \item \cite{baker_marinemammals_2010}
  \item \cite{baker_census_2010}
\end{itemize}

% Bibliography should be printed starting on a new page.
\clearpage
\printbibliography

%End of document
\end{document}
